% Generated by render.py; edit preamble.tex, citations.json and main.tikz
\documentclass[border=12pt]{standalone}
\def\pgfsysdriver{pgfsys-dvisvgm.def}
\usepackage[dvisvgm]{xcolor}
\usepackage{xeCJK}
\setCJKmainfont{HaranoAjiGothic-Regular.otf}[BoldFont=HaranoAjiGothic-Bold.otf]
\setCJKsansfont{HaranoAjiGothic-Regular.otf}[BoldFont=HaranoAjiGothic-Bold.otf]
\usepackage{amsmath,amssymb,tikz,lmodern}
\usetikzlibrary{arrows.meta,positioning,calc,fit,backgrounds}
\renewcommand{\familydefault}{\sfdefault}
\definecolor{ink}{HTML}{25322F}
\definecolor{green}{HTML}{246657}
\definecolor{tint}{HTML}{EEF4F0}
\definecolor{soft}{HTML}{C6D2CC}
\definecolor{muted}{HTML}{52615B}
\newif\ifJapanese
\newcommand{\Tx}[2]{\ifJapanese #1\else #2\fi}
\newcommand{\Eff}{\overline{\operatorname{Eff}}}
\newcommand{\heading}[3]{{\color{green}\bfseries #1\quad\Tx{#2}{#3}}\\[5pt]}
% A white rectangle makes the complete citation label clickable, including spaces.
\newcommand{\refurl}[2]{\special{dvisvgm:raw <a xlink:href="#1" target="_blank" rel="noopener noreferrer">}{\setlength{\fboxsep}{0pt}\colorbox{white}{\color{green}\underline{\strut #2}}}\special{dvisvgm:raw </a>}}
\tikzset{
 card/.style={draw=soft,fill=white,rounded corners=2pt,align=left,text=ink,font=\fontsize{10.5}{14}\selectfont,inner xsep=12pt,inner ysep=10pt,text width=14.4cm},
 half/.style={card,text width=6.45cm,minimum height=2.05cm},
 result/.style={card,draw=green,fill=tint,line width=.9pt},
 reason/.style={align=left,text=muted,font=\fontsize{9}{12}\selectfont,inner sep=0pt,text width=12.6cm},
 branchreason/.style={reason,text width=5.95cm},
 edge/.style={-{Stealth[length=2.4mm]},draw=green,line width=1pt},
 wire/.style={draw=green,line width=1pt},
}
% Keep the arrow in a dedicated left gutter; explanations never cover it.
\def\flow#1#2#3;{
 \coordinate (#2-from) at ($(#1.south)+(-6.9,0)$);
 \coordinate (#2-to) at ($(#2.north)+(-6.9,0)$);
 \draw[edge] (#2-from)--(#2-to);
 \node[reason,anchor=west] (#2-reason) at ($(#2-from)!.5!(#2-to)+(.55,0)$) {#3};
}

% Style adapted from the repository's Schnell diagrams; MIT attribution retained.

% Generated from citations.json; do not edit.
\newcommand{\Cite}[1]{\csname cite-#1\endcsname}
\expandafter\def\csname cite-main-result\endcsname{\shortstack[l]{\refurl{https://github.com/openai/math/blob/1581d4cbdd69a0faf34b2089614fe2a814316a6a/preprints/Global-Spherical-Shells-on-Minimal-Surfaces-of-Class-VII-September-24-2026/paper.pdf}{Theorem 1.1 · p. 1}}}
\expandafter\def\csname cite-reductions\endcsname{\shortstack[l]{\refurl{https://github.com/openai/math/blob/1581d4cbdd69a0faf34b2089614fe2a814316a6a/preprints/Global-Spherical-Shells-on-Minimal-Surfaces-of-Class-VII-September-24-2026/paper.pdf}{Proposition 2.1; Lemmas 2.2–2.5 · pp. 5–8}}}
\expandafter\def\csname cite-classes\endcsname{\shortstack[l]{\refurl{https://github.com/openai/math/blob/1581d4cbdd69a0faf34b2089614fe2a814316a6a/preprints/Global-Spherical-Shells-on-Minimal-Surfaces-of-Class-VII-September-24-2026/paper.pdf}{Lemma 3.1 · pp. 8–9}}}
\expandafter\def\csname cite-sections\endcsname{\shortstack[l]{\refurl{https://github.com/openai/math/blob/1581d4cbdd69a0faf34b2089614fe2a814316a6a/preprints/Global-Spherical-Shells-on-Minimal-Surfaces-of-Class-VII-September-24-2026/paper.pdf}{Propositions 4.8, 4.10–4.11 · pp. 17–20}}}
\expandafter\def\csname cite-residues\endcsname{\shortstack[l]{\refurl{https://github.com/openai/math/blob/1581d4cbdd69a0faf34b2089614fe2a814316a6a/preprints/Global-Spherical-Shells-on-Minimal-Surfaces-of-Class-VII-September-24-2026/paper.pdf}{Theorem 5.2; Propositions 5.3–5.4 · pp. 23–26}}}
\expandafter\def\csname cite-cylinder\endcsname{\shortstack[l]{\refurl{https://github.com/openai/math/blob/1581d4cbdd69a0faf34b2089614fe2a814316a6a/preprints/Global-Spherical-Shells-on-Minimal-Surfaces-of-Class-VII-September-24-2026/paper.pdf}{Proposition 6.3; Corollary 6.4 · pp. 30–32}}}
\expandafter\def\csname cite-compactness\endcsname{\shortstack[l]{\refurl{https://github.com/openai/math/blob/1581d4cbdd69a0faf34b2089614fe2a814316a6a/preprints/Global-Spherical-Shells-on-Minimal-Surfaces-of-Class-VII-September-24-2026/paper.pdf}{Lemma 6.5; Theorem 6.6 · pp. 32–34}}}
\expandafter\def\csname cite-count\endcsname{\shortstack[l]{\refurl{https://github.com/openai/math/blob/1581d4cbdd69a0faf34b2089614fe2a814316a6a/preprints/Global-Spherical-Shells-on-Minimal-Surfaces-of-Class-VII-September-24-2026/paper.pdf}{Lemmas 7.1–7.2; Proposition 7.3; Lemma 7.4 · pp. 34–37}}}
\expandafter\def\csname cite-finish\endcsname{\shortstack[l]{\refurl{https://github.com/openai/math/blob/1581d4cbdd69a0faf34b2089614fe2a814316a6a/preprints/Global-Spherical-Shells-on-Minimal-Surfaces-of-Class-VII-September-24-2026/paper.pdf}{Proof of Theorem 1.1 · p. 37}}}
\expandafter\def\csname cite-char\endcsname{\shortstack[l]{\refurl{https://github.com/openai/math/blob/1581d4cbdd69a0faf34b2089614fe2a814316a6a/preprints/Global-Spherical-Shells-on-Minimal-Surfaces-of-Class-VII-September-24-2026/paper.pdf}{Lemma 2.3; Proposition 4.1 · pp. 6, 11–12}}}
\expandafter\def\csname cite-index\endcsname{\shortstack[l]{\refurl{https://github.com/openai/math/blob/1581d4cbdd69a0faf34b2089614fe2a814316a6a/preprints/Global-Spherical-Shells-on-Minimal-Surfaces-of-Class-VII-September-24-2026/paper.pdf}{Propositions 4.10–4.11 · pp. 18–20}}}
\expandafter\def\csname cite-growth\endcsname{\shortstack[l]{\refurl{https://github.com/openai/math/blob/1581d4cbdd69a0faf34b2089614fe2a814316a6a/preprints/Global-Spherical-Shells-on-Minimal-Surfaces-of-Class-VII-September-24-2026/paper.pdf}{Lemma 5.1; equation (5.4) · pp. 21–23}}}
\expandafter\def\csname cite-curve-bound\endcsname{\shortstack[l]{\refurl{https://github.com/openai/math/blob/1581d4cbdd69a0faf34b2089614fe2a814316a6a/preprints/Global-Spherical-Shells-on-Minimal-Surfaces-of-Class-VII-September-24-2026/paper.pdf}{Proposition 7.3; Lemma 7.4 · pp. 36–37}}}
\expandafter\def\csname cite-enoki\endcsname{\shortstack[l]{\refurl{https://www.jstage.jst.go.jp/article/tmj1949/33/4/33_4_453/_pdf/-char/en}{[Eno81, Main Theorem; note (3), print p. 454 / PDF p. 2]}}}
\expandafter\def\csname cite-dot\endcsname{\shortstack[l]{\refurl{https://www.jstage.jst.go.jp/article/tmj/55/2/55_2_283/_pdf/-char/en}{[DOT03, Main Theorem; §1 (GSS definition), print pp. 283–284 / PDF pp.}\\\refurl{https://www.jstage.jst.go.jp/article/tmj/55/2/55_2_283/_pdf/-char/en}{1–2]}}}
\expandafter\def\csname cite-teleman\endcsname{\shortstack[l]{\refurl{https://annals.math.princeton.edu/wp-content/uploads/annals-v172-n3-p07-p.pdf\string##page=11}{[Tel10, Enoki surfaces, print p. 1758 / PDF p. 11]}}}
\expandafter\def\csname cite-basis\endcsname{\shortstack[l]{\refurl{https://annals.math.princeton.edu/wp-content/uploads/annals-v172-n3-p07-p.pdf\string##page=9}{[Tel10, §1.1 (canonical basis), print p. 1756 / PDF p. 9]}}}
\expandafter\def\csname cite-planar\endcsname{\shortstack[l]{\refurl{https://www.impan.pl/shop/en/publication/transaction/download/product/84985}{[Hem11, Theorem 1.1, print p. 77 / PDF p. 1]}}}
\expandafter\def\csname cite-oka-use\endcsname{\shortstack[l]{\refurl{https://github.com/openai/math/blob/1581d4cbdd69a0faf34b2089614fe2a814316a6a/preprints/Every-complex-K3-surface-is-Oka-September-23-2026/paper.pdf}{[OpenAI26k3, Corollary 10.2, p. 54]}}}

\Japanesefalse
\begin{document}
\begin{tikzpicture}
\node[result] (start) at (0,0) {
 \heading{1}{球殻判定で場合を分ける}{Separate the known shell case}
 $X$ \Tx{連結・極小・コンパクト複素曲面}{connected, minimal, compact complex surface},
 $b_1=1,\ \kappa=-\infty,\ b=b_2>0$.\\[4pt]
 \Tx{非零有効因子 $C$ に $C^2=0$}{An effective $C\ne0$ has $C^2=0$}
 $\ \Longrightarrow\ $\Tx{大域球殻}{global spherical shell}.\\[4pt]
 {\fontsize{9}{12}\selectfont\Cite{teleman}}\\[4pt]
 \Tx{以下はすべての非零有効因子に $C^2<0$ の枝。}{Continue below when every nonzero effective divisor has $C^2<0$.}
};
\node[card,below=1.8cm of start] (classes) {
 \heading{2}{切断して多数の相対類を保つ}{Cut and retain many relative classes}
 $Y\to X,\quad X_n=Y/\langle T^n\rangle,\quad d\ge nb-C_0$,\\[3pt]
 $e_i\in H_c^2(Y;\mathbf Z),\qquad e_i^2=K_Y\cdot e_i=-m_i<0$.
};
\flow{start}{classes}{\Tx{負定値格子の基底元の和を切断面上で自明化する}{Trivialize sums from the negative definite lattice at the seam}\\[3pt]\Cite{classes}};
\node[card,below=2cm of classes] (sections) {
 \heading{3}{指数1から零因子へ}{Pass from index one to a zero divisor}
 $\operatorname{ind}\bar\partial_{E_i}=1,\quad s_i\to1\ (h\to+\infty),\quad s_i\to0\ (h\to-\infty)$.\\[3pt]
 $D_i=(s_i=0),\quad [D_i]=j(e_i),\quad \sup_{D_i}h<\infty$.\\[3pt]
 \Tx{因子は局所有限。成分のコンパクト性はまだ未確定。}{The divisor is locally finite; component compactness is still to be proved.}
};
\flow{classes}{sections}{\Tx{重み付き消滅・正端一意性と切除を使う}{Use weighted vanishing, positive-tail uniqueness, and excision}\\[3pt]\Cite{index}};
\node[card,below=2.4cm of sections] (residue) {
 \heading{4}{スカラー留数で開正規化を制限する}{Restrict open normalizations by scalar residues}
 $\operatorname{Area}(D_i\cap Y_\ell)\le C_i e^{A\ell},\qquad \gamma>2A+10$.\\[3pt]
 $H^{2,1}_{-\varepsilon,\beta}(Y,\mathcal O)=0$
 \Tx{と留数独立性}{and residue independence}\\[3pt]
 $\Longrightarrow\quad R\subset\mathbf P^1,\quad
 R\text{ \Tx{非コンパクト}{noncompact}}\Longrightarrow R\simeq\mathbf C$.
};
\flow{sections}{residue}{\Tx{$L^2$ 微分を排除し、円筒評価で $dz/z$ も排除する}{Exclude $L^2$ differentials, then use the cylinder estimate on $dz/z$}\\[3pt]\Cite{residues}\\[3pt]\Cite{cylinder}};
\node[card,below=2.3cm of residue] (compact) {
 \heading{5}{同じ束・同じ重みで平面を排除する}{Exclude planes in one fixed weighted bundle}
 $B_k=T^kB\not\subset D_i,\quad \alpha_k=(f_k^*s_i)\,dz$.\\[3pt]
 $H=e^g\sigma=O(|z|^{-1}),\quad
 \operatorname{Mass}((f_k)_*\alpha_k;Y_\ell)\le C_k e^{A\ell}$.\\[3pt]
 \Tx{固定した有限次元}{One fixed finite-dimensional}
 $H^{2,1}_{-\varepsilon,\beta}(Y,E_i)$\\[3pt]\Tx{に無限独立族は入らない。}{Infinitely many independent residue classes are impossible.}\\[3pt]
 $\Longrightarrow\quad$\Tx{各成分はコンパクトで有理正規化を持つ}{every component is compact rational}.
};
\flow{residue}{compact}{\Tx{周期ゲージと2度目の円筒評価。$\varepsilon>0$, $\beta>\max(A,\varepsilon)$を固定}{Periodic gauge and a second cylinder estimate; fix $\varepsilon>0$, $\beta>\max(A,\varepsilon)$}\\[3pt]\Cite{compactness}};
\node[card,below=2cm of compact] (count) {
 \heading{6}{Laurent交点形式で曲線を数える}{Count curves by the Laurent intersection form}
 $\Lambda=\mathbf Q[t,t^{-1}],\quad P(t)=(\sum_k\langle e_i,S_!^ke_j\rangle t^k)_{ij}$.\\[3pt]
 $\det P(1)\ne0,\quad d\ge nb-C_0,\quad r<b\Longrightarrow nr<d\ (n\gg0)$.\\[3pt]
 \Tx{軌道数不足が非退化性に矛盾。曲線類の独立性も併用。}{Too few lift orbits contradict nonsingularity; also use curve-class independence.}\\[3pt]
 $\Longrightarrow\quad$\Tx{有理曲線がちょうど $b$ 本}{exactly $b$ rational curves}.
};
\flow{compact}{count}{\Tx{コンパクト試験類を局所有限因子と交差させる}{Pair compact test classes with locally finite divisors}\\[3pt]\Cite{count}};
\node[result,below=2.8cm of count] (shell) {
 \heading{7}{正則球殻と補集合条件を回収する}{Recover the holomorphic shell and its complement}
 $S^3\subset U\subset\mathbf C^2\setminus\{0\},\quad
 \varphi:U\overset{\sim}{\longrightarrow}\Sigma\subset X$.\\[3pt]
 $U,\Sigma$ \Tx{は開集合、$X\setminus\Sigma$ は連結}{are open, and $X\setminus\Sigma$ is connected}.
};
\flow{count}{shell}{\Tx{曲線判定の全仮定を満たし、球殻の定義まで使う}{Apply the curve criterion, including the full shell definition}\\[3pt]\Cite{dot}\\[3pt]\Cite{finish}};

\end{tikzpicture}
\end{document}
